\section{ПРОЕКТИРОВАНИЕ СИСТЕМЫ ИНДЕКСАЦИИ ОНЧЕЙН-ДАННЫХ}

После анализа предметной области следующим этапом является проектирование
системы, обеспечивающей устойчивую индексацию, хранение и предоставление
ончейн-данных. На данном этапе необходимо формализовать требования,
определить архитектурный подход, спроектировать состав программных модулей,
модель хранения данных и интерфейсы взаимодействия с внешними потребителями.

Проектирование системы Bitcoin Blockchain Indexer выполнялось с учетом
ограничений предметной области и требований к эксплуатационной надежности.
В качестве исходных принципов использовались модульность, управляемость
жизненного цикла заданий индексирования, транзакционная согласованность
данных и пригодность к контейнерному развертыванию
\cite{gost_arch, gost_quality, kleppmann2017designing, richards2020fundamentals, postgres_docs}. Тем самым проектирование
рассматривается не как произвольный выбор реализации, а как процесс
обоснования архитектурных решений в соответствии с выявленными требованиями.
